\section{Parameter selection without circularity}

Fix \(a,\nu\) first. Thus \(\rho\) and \(\delta=\log d\) are fixed.

Choose positive rational \(\theta,A,B,C\) satisfying
\[
0<\theta<A<B<1,\quad
\nu(A-\theta)>1-\theta,\quad
C>1,\quad CB<1,\quad B<C\theta<1.
\tag{16}
\]
For existence, choose rational \(b\in(1/2,1-1/\nu)\), then a small positive rational \(s\), and put \(\theta=1-s,A=1-bs\). For small \(s\), the gap in (16) is positive and \(A^2<\theta\). Choose \(C\in(A/\theta,1/A)\) and then \(B\in(A,\min(1/C,C\theta))\).

Choose a small rational \(\eta>0\) so that
\[
g=\nu(A(1-\eta)-\theta)-(1-\theta)>0,
\qquad \epsilon=\min(g/2,1/2).
\]
Take \(F_0>2/\theta\) with \(\nu/F_0<\epsilon/3\).

Now increase \(m\), setting
\[
K=\lfloor C^m\rfloor,\quad w_0=B^{-m},
\quad v_0=2K\theta^mw_0.
\]
The ambient interpolation mass is exactly \(1/2\). Also
\[
\frac K{w_0}\le(CB)^m\to0,\quad
v_0\sim2(C\theta/B)^m\to\infty,\quad
\frac m{v_0}\to0,
\]
while (15) becomes
\[
c_\infty=\frac{c\eta^2(B/A)^m}{2(m+1)}\to\infty.
\]
Choose one finite \(m\ge1\) for which
\[
\frac{\Lambda F_0m+2\log2}{v_0}
+\frac{(\rho+\delta)K}{w_0}<\epsilon/3,
\qquad c_\infty>2.
\tag{17}
\]
Fix these parameters.

Since \(K\theta^m<1\), choose small rational \(\sigma>0\) satisfying
\[
(1+3\sigma)^{m+1}/2<1,\quad
(1+3\sigma)^mK\theta^m<1,\quad
(1+\sigma)\theta<1.
\]
The middle condition is retained in the formal parameter package; the distinct-\(Y\) interpolation proof uses the ambient mass inequality.

Only now choose the finitely many approximations. Arbitrarily large denominators let us make all \(w_i\) large and successively separated, satisfying (3), as well as
\[
\Lambda\sum_i\frac1{w_i}
+\frac{\theta+\log4+\log(2K)+\nu+\log(2\rho K)}{w_*}
<\epsilon/3.
\tag{18}
\]
There are finitely many product comparisons. Since \(v_i=w_i/\theta\) for positive indices, common factors cancel and each remaining comparison is met by increasing the next denominator sufficiently. The weight-separation requirement and the uniform rectangular cutoff depend on the fixed weights and \(\sigma\), rather than the center locations. The eventual interpolation threshold may depend on the fully selected centers; these are fixed before \(H\) tends to infinity.

Equations (12), (17), and (18) give
\[
E_{\rm ar}+E_{\rm an}<\epsilon<g,\qquad
c_\infty>1+E_{\rm ar}+E_{\rm an}.
\tag{19}
\]
This order absorbs both costs introduced by a general rational base: \(\log d\) in the arithmetic estimate and \(\rho\) in the analytic radius.

\section{Comparing the same nonzero determinant}

For sufficiently large cofinal \(H=nR\), interpolation supplies the minor, so (10) and (14) both apply.

If the maximum in (14) is the collision alternative, compatibility requires
\[
c_H\le1-b_H+E_{\rm ar}+E_{\rm an}+o(1)
\le1+E_{\rm ar}+E_{\rm an}+o(1).
\]
This contradicts the strict limiting margin in (19).

For the other alternative compatibility requires
\[
\nu(A(1-\eta)-b_H)-(1-b_H)
\le E_{\rm ar}+E_{\rm an}+o(1).
\]
The left side is at least \(g\): it decreases with \(b_H\), and \(b_H\le\theta\). This contradicts the other margin in (19).

Both alternatives fail. Hence arbitrarily large approximations of the stated quality cannot exist, proving (1). The irrationality and Dirichlet argument from §1 completes the exponent-two conclusion.
